% РЕФЕРАТ (стр. 3 записки — после титула и задания, перед ведомостью и оглавлением)
% Не выносится в оглавление. Должен умещаться строго на одну страницу.
% Счётчики рис./табл./источников и номер последней страницы заключения
% подставляются автоматически (\totfig, \tottab, \totref, \pageref{lastmainpage}).
% Число приложений — вручную.
\clearpage
\setcounter{page}{3}
\thispagestyle{empty}

{\centering\textbf{РЕФЕРАТ}\par}
\vspace{12pt}

ПРИЁМНАЯ КАМПАНИЯ, АВТОМАТИЗАЦИЯ, АЛГОРИТМ ОТБОРА, ВЕБ-ПРИЛОЖЕНИЕ, ВАЛИДАЦИЯ ДАННЫХ

Объектом разработки является программное обеспечение для компьютеризации ввода и мониторинга информации об абитуриентах при поступлении в Белорусский национальный технический университет, включающее алгоритм автоматического распределения абитуриентов по конкурсным спискам и категориям приёма.

Целью проекта стало проектирование и реализация централизованной информационной системы, способной автоматически вести учёт данных абитуриентов, ранжировать их по настраиваемым критериям и распределять места с учётом планов приёма, квот и приоритетов категорий, а также контролировать достоверность вносимых данных.

В процессе проектирования выполнены анализ предметной области, разработка модели данных и архитектуры приложения, реализация алгоритма отбора по принципу очереди с вытеснением, а также построение подсистем аудита, двойного контроля данных и отложенного удаления записей.
Проведено функциональное тестирование системы.

Практическая значимость заключается в замене разрозненных средств учёта~--~электронных таблиц и ручного пересчёта рейтингов~--~единой централизованной системой, что повышает достоверность и прозрачность распределения мест, снижает трудозатраты приёмной комиссии и исключает ошибки ручного ранжирования.

Область применения~--~приёмные комиссии университетов и иных учреждений образования, осуществляющих конкурсный отбор абитуриентов по совокупности оценочных показателей.

В рамках проекта апробированы гибкая настройка критериев и групп ранжирования, разграничение прав по ролям и журналирование всех действий; архитектура <<React~--~ASP.NET~Core~--~PostgreSQL>> обеспечивает масштабируемость и сопровождаемость решения.

Студент-дипломник подтверждает, что приведённый в дипломном проекте расчётно-аналитический материал объективно отражает состояние разрабатываемого объекта, все заимствованные из литературных и других источников теоретические и методологические положения сопровождаются ссылками на их авторов.

Дипломный проект: \pageref{lastmainpage}~с., \totfig~рис., \tottab~табл., \totref~источников, 1~прил.
\clearpage
